\section{Introduction}
\label{sec:intro}

3D scene reconstruction and novel view synthesis aim to recover
photorealistic scene representations from multi-view images. Neural Radiance
Fields (NeRF)~\cite{nerf} demonstrated that implicit radiance-field
representations can produce high-fidelity view synthesis from posed images,
while 3D Gaussian Splatting (3DGS)~\cite{kerbl20233d} made this line of
research more practical for real-time rendering by representing scenes with
explicit 3D Gaussians. Thanks to its strong reconstruction quality, fast
rendering speed and scene editability, 3DGS has quickly become a new paradigm
for scene reconstruction and novel view synthesis~\cite{survey3dgs}.

However, the success of vanilla NeRF and 3DGS still depends heavily on the
reliability of image supervision. This assumption breaks down severely in
low-light conditions. Under adverse illumination, image data often suffer
from severe noise amplification, color distortion, loss of local detail, and
photometric inconsistency across viewpoints~\cite{liu2025realx3d,zhou2025lita}.
More importantly, these degradations are not uniformly distributed: some
regions still preserve usable structure and appearance, while others become
completely unreliable. Luminance-GS~\cite{cui2025luminance} explicitly points
out that low-light and exposure variations introduce photometric
inconsistencies across views for both NeRF- and 3DGS-based pipelines, while
LITA-GS~\cite{zhou2025lita} further highlights the lack of reliable
references, severe information loss, and insufficient SfM points under low
illumination. As a result, low-light 3D reconstruction could be seen as a
problem of heterogeneous reliability: the supervision signal and the
available priors are unevenly trustworthy across pixels, regions and
optimization stages.

Prior work has made important progress toward illumination-robust
reconstruction. AlethNeRF~\cite{cui2024aleth} introduced concealing fields to
explain low-light degradation and enabled unsupervised recovery of well-lit
renderings from dark observations. I2NeRF~\cite{liu2025i2nerf} formulated
media degradation through a radiative model that unifies emission,
absorption, and scattering under the Beer-Lambert law. More recently,
Luminance-GS~\cite{cui2025luminance} adapted 3DGS to challenging lighting
conditions through per-view color mapping and view-adaptive curve adjustment,
and LITA-GS~\cite{zhou2025lita} introduced illumination-invariant physical
priors together with a progressive denoising strategy for reference-free
3DGS under adverse lighting. These methods significantly improve robustness
under illumination changes, but they mostly emphasize degradation modeling,
photometric adaptation, or illumination-invariant representations.

In low-light 3D reconstruction, supervisory cues are not uniformly reliable
across either image regions or training stages. Motivated by this
observation, we present Reliability-Aware Staged Low-Light 3DGS decomposing
the problem into three stages: geometry bootstrap, reliability-aware geometry
refinement, and geometry-frozen appearance refinement. In the first stage, we build a
stable scene scaffold with weak depth only used as a conservative trend
prior~\cite{marigold}. In the second stage, we warm-start from this scaffold
and continue to reuse sparse geometry from Structure-from-Motion
(SfM)~\cite{sfm_revisited} during training as quality-aware soft support.
In the final stage, we freeze geometry and perform geometry-frozen appearance
refinement, where luminance restoration and chromatic correction are explicitly
decoupled in YCbCr space to reflect their different reliability in low-light
images. This design reflects our central claim: low-light 3DGS should not
treat all supervision equally, but should exploit each cue only when and
where it is reliable enough to help.

Our contributions can be summarized as follows:

\begin{enumerate}
    \item We design a reliability-aware pipeline, formulating low-light 3DGS
    as a reliability-aware optimization problem, where supervisory cues and
    priors have heterogeneous trustworthiness across regions and stages.

    \item We introduce a density/quality-aware soft sparse anchoring
    strategy, where fixed-pose COLMAP sparse points are used not only for
    initialization but also as weak geometric support during refinement.

    \item We propose a decoupled appearance refinement strategy in YCbCr
    space, where luminance and chroma are supervised by different
    targets with region confidence awareness.
\end{enumerate}
